
% --------------------------------------------------------------------------------------------------
% Kommando Beispiele
% --------------------------------------------------------------------------------------------------

% Referenz auf z.B. Kapitel
\chapter{Introduction}
\label{c:introduction}

%Here how to refer a chapter \ref{c:introduction}...

% Literaturhinweis Reference example \cite{arendt2010henshin}


% Theorem Style
%\begin{definition}\label{def:emc} $e=mc^2$ \end{definition}
%\begin{lemma} $e=mc^2$ \end{lemma}
%\begin{proof} $e=mc^2$ \end{proof}
%\begin{requirement} $e=mc^2$ \end{requirement}
%\begin{example} $e=mc^2$ \end{example}

% Here how to refer a Definition: In Definition \ref{def:emc}...

% Noch zu Erledigen Hinweis
% \todo{HERE A TODO!}

% Randnotiz
%\notiz{SOME NOTE}

% Struktur
\section{Motivation and Purposen}
\subsection{SubSection}
\subsubsection{SubSubSection}
\paragraph{Pragraph}
\subparagraph{SubPragraph}

\section{Limitations and AI}
How AI can help. Input from Podcast SRF. Terence Tao thinks in 3 years more helpful, the other from Lean, thinks in 3 years Math is "solved", like chess, where humans only play a minor role. Lean will soon integrate some AI to guide the Mathematician with the proof.

\subsection{Formalizing Fermat's Last Theorem}
4.9.2026: Claude worked largely autonomously over 11 days to write the proof in the Lean programming language. \href{https://www.anthropic.com/research/formalizing-fermats-last-theorem}{Link anthropic}

\section{Hello World!}
As Lean is a pure functional language, printing "Hello World!" to the console isn't pure at all. But as a computer scientist, you can imagine what's next. \\ 
Installing Lean is stright forward and best using VS Code and the Lean 4 extension form leanpower. Accept installing Elan, the version manager for Lean, as VS Code promts for it. This can take some minutes. Restart VS Code, open a new file and store it with the extenson \texttt{lean}. Now you are ready to eval your first experssion in the lean file and to check the version of lean in the terminal by typing:
\begin{verbatim}
    > lean --version
    Lean (version 4.33.1, ...)
\end{verbatim}

\section{Concepts}

Several so called \textit{tactics} exist, to prove a theorem:
\begin{itemize}
    \item rw - rewrite, substitutions
    \item ring - algebraic manipulations
\end{itemize}

You are only allowed to use theorems, if the preconditions are verified too. E.g. using the product rule
$$ (u,v)' = uv' + u'v$$
requires to prove the derivability of $u$ and $v$, annoying. Or using 
$$\sum^{N}_{n=0} = \frac{1}{1-x}$$
is only allowed if you show $|x|<1$ and $N \geq 0$. \\
This makes Lean rigorous.Ref https://www.youtube.com/watch?v=bldhmF1Isk4.

\paragraph{Propsitions, Proofs and Types}
A proposition can be true or false. Like $1+1=2$ or "it's raining tomorrow". We only concern the former ones, the purely mathematical propositions relying only on definitions.\\

A proof on the other hand, is a confirmation or argument, that a proposition is true and is based on 
\begin{itemize}
    \item the mathematical definitions and
    \item logical argumentation.
\end{itemize}

Mathematicians understand the proof and leave out details. Lean fills in missing steps and checks them mechanically to eliminate possible faults in the proof.\\

Beside functions, also propositions are types. The type of a propositions specifys the evidence that it is true. To prove a proposition, Lean checks those evidences. \textit{In Lean, proving $P$ means constructing a term (or value) of type $P$. Let's look at the following exampele where \texttt{rfl} is a constructor that produces such a term for certain equality propositions}.\\

For the proposition $1+1=2$ the constructor \texttt{rfl} produces a value (term) of type $P$ as the proof of $P$. If the propositon is False, the constructor fails and the proof cannot be constructed.

\begin{verbatim}
    def foo : 1 + 1 = 2 := rfl  
\end{verbatim}
telling \texttt{foo} is of type $1+1=2$, and \texttt{rfl} constructs the proof if the two sides are equal. As Lean can evaluate $1+1$ on the left side, the equality is given. The constructor \texttt{rfl} cannot construct the proposition $1+1=55$. 


